% =====================================================================================================
\section{Machine-learning methods}\label{sec:ml}
% =====================================================================================================

\subsection{Problem formulation and software}\label{sec:mlsetup}
Each model is a binary classifier on observations: it maps a feature vector $\vect x_i$ to a BH score
$s(\vect x_i)\in[0,1]$. The question is about \emph{sources}, so every score is produced out of fold (no model sees the
source it scores), and observation scores are averaged per source (\cref{eq:sourcescore}). All models are scikit-learn
1.8.0 estimators \citep{Pedregosa2011} (\file{requirements-lock.txt}) with the settings in \file{config.py}. They are wrapped in the pipelines of
\file{07_train_eval.py} and \file{v3lib.py}, in \code{V4Model} (\file{v4lib.py}, the eight-algorithm search) and in
\code{C2Model} (\file{v7lib.py}, MAXI). \Cref{tab:models} lists them.

\begin{table}[htbp]
\centering\small
\caption{Models and their roles. P = model of a planned primary comparison; B = baseline; D/E = descriptive or exploratory.}\label{tab:models}
\begin{tabularx}{\textwidth}{@{}>{\raggedright\arraybackslash}p{3.6cm}>{\raggedright\arraybackslash}p{3.3cm}>{\raggedright\arraybackslash}X c@{}}
\toprule
Model & Rounds & Code & Role \\
\midrule
Logistic regression (LR) & v1--v13 & \file{07}, \code{v3lib.MODELS}, \code{V4Model}, \code{C2Model} & P, B \\
Random forest (RF) & v1--v13 & as LR & P (v8d, v12b, v13a), D \\
Dummy (class prior) & v1, v2 & \file{07_train_eval.py} & B \\
ExtraTrees, HistGB, QDA, MLP & v4a & \code{V4Model} & E \\
RBF SVM & v4a, v7c2; v7c1 & \code{V4Model}, \code{C2Model}, \file{37} & E, reproduction \\
kNN classifier & v4a, v7c2, v8d; v7c1 & \code{V4Model}, \code{C2Model}, \file{37} & E, reproduction \\
Nested automatic selection & v4a & \code{v4lib.nested_fold} & P (v4) \\
CCTLR & v6c & \code{v6lib.CCTLR} & D \\
kNN regression (residualisation) & v6d, v9b--v13b & \code{v6lib.source_residuals} & inside the permutation tests \\
Frozen $\nu^\ast$ rule & v11, v13 & \file{57}, \file{63} & D \\
\bottomrule
\end{tabularx}
\end{table}

% -----------------------------------------------------------------------------------------------------
\subsection{Shared components}\label{sec:shared}
\paragraph{Standardisation.} Scale-sensitive models (LR, SVM, kNN, QDA, MLP, the colour part of CCTLR, the kNN regression)
standardise each feature with the training rows $\mathcal T$ of the current fit:
\begin{equation}
z_{i\ell}=\frac{x_{i\ell}-\mu_\ell}{\sigma_\ell},\qquad\mu_\ell=\frac1{|\mathcal T|}\sum_{i\in\mathcal T}x_{i\ell},\qquad
\sigma_\ell^2=\frac1{|\mathcal T|}\sum_{i\in\mathcal T}(x_{i\ell}-\mu_\ell)^2\quad(\sigma_\ell=0\Rightarrow1),\label{eq:scaler}
\end{equation}
and apply the same $(\mu_\ell,\sigma_\ell)$ to the test rows (\code{StandardScaler} inside the pipeline). Tree models use
the raw features.

\paragraph{Class-balanced weights.}\label{sec:weights} The main sample has 7 BH and 24 NS sources. With $n$ training rows,
$n_c$ of them in class $c$, \code{class_weight='balanced'} gives each row
\begin{equation}
w^\mathrm{bal}_i=\frac{n}{2\,n_{y_i}},\qquad\sum_{i:y_i=c}w^\mathrm{bal}_i=\frac n2.\label{eq:balanced}
\end{equation}
QDA takes no weights and uses equal priors instead. kNN corrects its output to equal priors (\cref{eq:knnprior}), and the MLP
receives $w^\mathrm{bal}$ as sample weights.

\paragraph{Source-equal weights (v4a, v4b1, v5 S3, v6 S3).} When sources contribute very different numbers of observations
(6 to 1329 per source in v4b1), balanced weights let the most-observed sources dominate. With $S_c$ training sources in
class $c$ and $n_k$ rows of source $k$ (\code{v4lib.source_equal_weights}):
\begin{equation}
w^\mathrm{src}_i=\frac{n}{2\,S_{y_i}\,n_{g_i}},\qquad\sum_{i:g_i=k}w^\mathrm{src}_i=\frac{n}{2S_c}\ (k\in c),\qquad\sum_{i:y_i=c}w^\mathrm{src}_i=\frac n2.\label{eq:sourceequal}
\end{equation}
Every source of a class then has the same total weight, and both classes sum to $n/2$, the scale of the balanced weights.
\code{class_weight} is set to \code{None} when these weights are passed. The first implementation normalised each class to 1
instead of $n/2$, which acted as a much stronger LR penalty. It was fixed after the v4a results and the affected parts were
re-run; no conclusion changed (decision log).

\paragraph{Scores.} \code{predict_proba[:,1]} for LR, RF, ExtraTrees, HistGB and QDA; the prior-corrected vote fraction for
kNN; the mean over seeds for the MLP; the logistic function of the decision value for SVM and CCTLR. Scores rank observations
and are thresholded at 0.5; they are not calibrated probabilities.

% -----------------------------------------------------------------------------------------------------
\subsection{Logistic regression}\label{sec:lr}
LR is the main model: H-LR (two colours) is the reference that every later model is compared against, and H+T-LR is the
primary model of v5, v6a, v8a/b, v9a, v10a and v11a. On standardised inputs, $s(\vect x)=\sigma(\vect w^\top\vect z+b)$
with $\sigma(u)=1/(1+e^{-u})$, and
\begin{equation}
(\hat{\vect w},\hat b)=\arg\min_{\vect w,b}\ \frac12\|\vect w\|_2^2+C\sum_{i\in\mathcal T}\omega_i\bigl[-y_i\log s(\vect x_i)-(1-y_i)\log(1-s(\vect x_i))\bigr],\label{eq:lrobj}
\end{equation}
with $C=1$, L-BFGS, at most 5000 iterations, an unpenalised intercept and $\omega_i=w^\mathrm{bal}_i$ (or $w^\mathrm{src}_i$).
With a few hundred rows and 2--6 features the penalty is weak. LR assumes log-odds that are linear in the standardised
features, so a timing feature has one slope everywhere in the colour plane. This is the motivation for CCTLR
(\cref{sec:cctlr}).

\subsection{Random forest}\label{sec:rf}
RF is the second reference model in every round, and the primary model of v8d (MAXI), v12b and v13a. It averages $B=500$
trees: $s(\vect x)=\frac1B\sum_b\hat p_b(\vect x)$, where $\hat p_b$ is the weighted BH fraction in the leaf of tree $b$. Each
tree is grown on a bootstrap sample (row weights $\omega_i$ times the bootstrap counts) with CART splits that maximise the
weighted decrease of the Gini impurity $G=1-\sum_c\pi_c^2$. At each node $m=\lfloor\sqrt p\rfloor$ features are candidates,
so with the two colours of H every split uses one randomly chosen colour, and with $p=6$ two. Trees grow until nodes are
pure or a child would have fewer than 2 rows (\code{min_samples_leaf=2}). Balanced class weights are computed once per fold,
and \code{random_state=42}. RF can follow non-monotone class regions, which matters for MAXI and for timing--colour
interactions. With fewer than 10 BH sources, however, a leaf can memorise one source.

% -----------------------------------------------------------------------------------------------------
\subsection{Further algorithms of the v4a search}\label{sec:otheralgs}
\begin{itemize}
\item \textbf{Dummy} (\code{strategy='prior'}): returns the training BH fraction and predicts the majority class. Under LOSO
the majority depends on the held-out source. In v1 (4+4 sources), holding out a BH leaves an NS majority and vice versa, so
the Dummy balanced accuracy is 0, not 0.5.
\item \textbf{ExtraTrees}: like RF but without bootstrap, and with one random threshold per candidate feature at each node.
\item \textbf{HistGradientBoosting}: an additive model on the log-odds scale; each of 200 trees is fitted to the weighted
gradients and Hessians of the log loss, with features pre-binned into at most 255 quantile bins. Learning rate 0.1, 15
leaves, at least 20 rows per leaf, no L2 term and no early stopping.
\item \textbf{RBF SVM}: \code{SVC} with balanced class weights and $K(\vect z,\vect z')=e^{-\gamma\|\vect z-\vect z'\|^2}$.
\code{gamma='scale'} gives $\gamma=1/(p\operatorname{Var}\vect Z)$; the v4a grid multiplies it by $\kappa\in\{0.3,1,3\}$.
The score is $\sigma(f(\vect z))$ of the decision value $f$. Platt scaling is not used, because it runs an internal random
cross-validation.
\item \textbf{kNN classifier}: the BH fraction $\hat p$ among the $k$ nearest standardised training rows (Euclidean,
uniform weights), corrected to equal priors with the training BH fraction $\pi$:
\begin{equation}
s(\vect x)=\frac{\hat p/\pi}{\hat p/\pi+(1-\hat p)/(1-\pi)}.\label{eq:knnprior}
\end{equation}
\item \textbf{QDA}: Gaussian classes with regularised covariances $\hat\Sigma_c=(1-r)S_c+rI$ ($r=0.1$) and equal priors.
\item \textbf{MLP}: one hidden layer of 32 ReLU units, Adam, $\alpha=10^{-3}$, balanced sample weights. Early stopping
holds out a stratified 20\% validation split and restores the best weights. Because a small MLP is unstable, the score is the
mean over seeds $\{0,\dots,4\}$ (outer fits) or $\{0,1,2\}$ (inner fits).
\end{itemize}
A planned 1D convolutional network on representation B was skipped, as declared in the plan, because \code{torch} was not
installed. Principal component analysis was never used: ``PCA'' in the code always means the Proportional Counter Array.

% -----------------------------------------------------------------------------------------------------
\subsection{Hyperparameters}\label{sec:hyper}
Except in the nested selection (\cref{sec:nested}) and the MAXI kNN and SVM, all hyperparameters are fixed in advance.
\Cref{tab:hyper} collects them.

{\small\setlength{\tabcolsep}{3pt}
\begin{longtable}{@{}>{\raggedright\arraybackslash}p{2.7cm}>{\raggedright\arraybackslash}p{6.3cm}>{\raggedright\arraybackslash}p{6.6cm}@{}}
\caption{Hyperparameters (\file{config.py}: \code{LR_PARAMS}, \code{RF_PARAMS}, \code{V4_FIXED}, \code{V4_GRID},
\code{V7_C2_*}, \code{V7_DEBEURS_*}).}\label{tab:hyper}\\
\toprule
Model & Fixed setting & Grid (where tuned) \\
\midrule\endfirsthead
\toprule Model & Fixed setting & Grid (where tuned) \\ \midrule\endhead
\bottomrule\endfoot
LR & $C=1$, L2, lbfgs, 5000 iterations, balanced & v4a: $C\in\{0.01,0.1,1,10\}$ \\
RF & 500 trees, \code{max_features='sqrt'}, leaf $\ge2$, balanced, seed 42 & v4a: leaf $\in\{1,2,5,10\}$; 200 trees inner, 500 outer \\
ExtraTrees & as RF & v4a: leaf $\in\{1,2,5,10\}$ \\
HistGB & learning rate 0.1, 200 iterations, 15 leaves, leaf $\ge20$, $\lambda=0$, no early stop, balanced, seed 42 & v4a: rate $\in\{0.05,0.1\}\times$ leaves $\in\{7,15\}$ \\
SVM (RBF) & $C=1$, $\gamma$ = \code{'scale'}, balanced & v4a: $C\in\{0.1,1,10\}\times\kappa\in\{0.3,1,3\}$; v7c2: $C\in\{0.1,1,10\}\times\gamma\in\{0.1,0.585,3\}$ \\
kNN & $k=15$, uniform, prior-corrected & v4a: $\{5,15,31\}$; v7c2, v8d: $\{5,15,24,35\}$ \\
QDA & $r=0.1$, priors $(0.5,0.5)$ & v4a: $r\in\{0.01,0.1,0.5\}$ \\
MLP & $(32,)$, $\alpha=10^{-3}$, early stopping (20\%), 500 epochs, seeds 0--4 & v4a: $\{(16,),(32,16)\}\times\alpha\in\{10^{-4},10^{-2}\}$ \\
CCTLR & Student $t$ with 4 d.o.f., scale floor 0.01, ridge $10^{-9}$ & none \\
kNN regression & $K=25$, uniform, neighbours from other sources & none \\
de Beurs reproduction & kNN $k=24$ on raw colours; SVM $C=0.655298$, $\gamma=0.585144$, Platt, one-vs-one, 3 classes & none (taken from the paper) \\
\end{longtable}}

% -----------------------------------------------------------------------------------------------------
\subsection{Model selection by nested leave-one-source-out}\label{sec:nested}
When anything is selected (algorithm, representation, hyperparameter or threshold), the selection runs inside each outer
training set, by an inner LOSO over its sources. The outer test source plays no part (\cref{fig:nested}).

\begin{figure}[htbp]
\centering
\begin{tikzpicture}[x=1cm, y=1cm,
  st/.style={fc model, text width=3.2cm, minimum height=15mm}]
  \node[fc data, text width=3.2cm, minimum height=15mm] (o) at (1.8,0) {outer fold $k$: 30 training sources; source $k$ held out};
  \node[st] (i) at (6.0,0) {for every candidate $c$: inner LOSO over the 30 sources $\Rightarrow$ inner source AUC$_c$, BA$_c$, threshold $t_c$};
  \node[st] (s) at (10.2,0) {choose $c^\ast$: highest inner AUC; ties by inner BA, then fixed order};
  \node[st, draw=red!70!black, fill=yellow!15] (f) at (14.4,0) {refit $c^\ast$ on the 30 sources (500 trees); score source $k$; store $\hat s-t^\ast+0.5$};
  \foreach \a/\b in {o/i,i/s,s/f} \draw[fc arr] (\a) -- (\b);
  \draw[fc arr, red!70!black] (o.south) -- ++(0,-0.55) -| node[fc note, pos=0.25, below] {the held-out source enters only here} (f.south);
\end{tikzpicture}
\caption{Nested LOSO (v4a). Every candidate is judged only on the outer training sources, so the cost of choosing the best
of many is included in the outer estimate.}\label{fig:nested}
\end{figure}

\paragraph{v4a (\code{v4lib.nested_fold}, \file{16_v4a_algorithms.py}).} Four representations (H, HI, B, A) are crossed with
eight algorithms and their grids (35 settings), giving 140 configurations. In part 1d, nesting runs per algorithm and
representation; in part 1e, over all 140 configurations. The inner stage uses 200 trees and MLP seeds $\{0,1,2\}$, the outer
refit 500 trees and seeds $\{0,\dots,4\}$.

\paragraph{Threshold choice (\code{v4lib.choose_threshold}).} With $u_1<\dots<u_q$ the distinct inner source scores, the
candidates are $u_1-10^{-9}$, the midpoints $(u_r+u_{r+1})/2$ and $u_q+10^{-9}$. Among those maximising the inner source
balanced accuracy, the one closest to 0.5 is taken:
\begin{equation}
t^\ast=\arg\min_{t\in\mathcal C^\ast}|t-0.5|,\qquad\mathcal C^\ast=\{t:\BA_\mathrm{src}(t)=\max_{t'}\BA_\mathrm{src}(t')\}.\label{eq:choosethreshold}
\end{equation}
The outer scores are stored shifted, $\hat s-t^\ast+0.5$, so that the common 0.5 rule applies.

\paragraph{MAXI (v7c2, v8d; \code{v7lib.c2_outer}).} kNN's $k$ is chosen by inner LOSO source AUC, and the SVM's $(C,\gamma)$
by inner \code{GroupKFold(5)} over sources (a pre-run change to save computing time). The first configuration with a strictly
higher AUC wins. Each training source contributes at most 200 randomly chosen days (seed 42), and the test source is scored
on all its days.

% -----------------------------------------------------------------------------------------------------
\subsection{Colour-conditional timing likelihood ratio (CCTLR, v6c)}\label{sec:cctlr}
In LR a timing feature has one slope everywhere, although timing is expected to help only where the colours overlap.
CCTLR adds to the H-LR logit a log-likelihood ratio of the timing features \emph{given} the colours (\code{v6lib.CCTLR}):
\begin{equation}
\ell(\vect x)=\vect w^\top\vect h+b+\sum_{k=1}^3\Bigl[\log t_4\bigl(T_k;\mu_{1k}(\vect h),\tau_{1k}\bigr)-\log t_4\bigl(T_k;\mu_{0k}(\vect h),\tau_{0k}\bigr)\Bigr],\qquad s=\sigma(\ell),\label{eq:cctlr}
\end{equation}
where $t_4$ is the Student-$t$ density with 4 degrees of freedom (robust to outliers such as bursts or dips) and
$\mu_{ck}(\vect h)=\beta_{ck0}+\beta_{ck1}h_1+\beta_{ck2}h_2$. For each class and timing band, $\vect\beta_{ck}$ is a weighted
least-squares fit on that class's rows with finite timing. The weights are $v_i=1/(S_cn_{g_i})$, so every source counts
equally, and a ridge of $10^{-9}$ is added. The scale $\tau_{ck}$ is the weighted residual standard deviation, floored at
0.01. If any timing value of a row is missing, the timing term is zero and the score equals the H-LR score. CCTLR has no
tunable parameter.

% -----------------------------------------------------------------------------------------------------
\subsection{kNN regression for colour residualisation}\label{sec:knnreg}
Inside every permutation test (\cref{sec:permutation}), the part of a timing feature $F$ (e.g.\ $\log_{10}\nu_c$) that is
predictable from colour is removed \emph{without labels} (\code{v6lib.source_residuals}). With standardised colours,
\begin{equation}
\hat F_i=\frac1K\sum_{j\in\mathcal N_K(\vect h_i;\ g_j\ne g_i)}F_j,\qquad r_i=F_i-\hat F_i,\qquad
\bar r_k=\frac1{n_k}\sum_{i:g_i=k}r_i,\label{eq:knnreg}
\end{equation}
with $K=25$ (capped at the number of available rows). The neighbours of an observation come only from \emph{other} sources,
so a source cannot explain itself. Rows with missing $F$ are skipped. If BHs cluster in one colour region, part of the class
difference is absorbed into $\hat F$, which makes the test conservative.

\subsection{Frozen characteristic-frequency rule (v11, v13)}\label{sec:nurule}
A simple rule frozen before the new data were read: a hard-like observation is called BH if $\nu_c<\nu^\ast$, and a source is
called BH if at least half of its new hard-like observations are. The threshold is the geometric mean of the class medians
in the previous data,
\begin{equation}
\nu^\ast=\sqrt{\median_\mathrm{BH}\nu_c\cdot\median_\mathrm{NS}\nu_c},\label{eq:nustar}
\end{equation}
giving 1.3007\,Hz from the v10 data (used in v11) and 1.3361\,Hz from the corrected v12 data (used in v13).
The values were written to \file{logs/v11/nu_star_frozen.txt} and \file{logs/v13/nu_star_frozen.txt} before the new
observations were read.

% -----------------------------------------------------------------------------------------------------
\subsection{Training designs}\label{sec:designs}
The same models are trained in five ways, depending on the question (\cref{tab:designs}).

\begin{table}[htbp]
\centering\small
\caption{Training designs. ``Out-of-source'': each source is scored by a model fitted on the reference data without that source.}\label{tab:designs}
\begin{tabularx}{\textwidth}{@{}>{\raggedright\arraybackslash}p{3.2cm}>{\raggedright\arraybackslash}X>{\raggedright\arraybackslash}p{4.2cm}@{}}
\toprule
Design & What is fitted and what is scored & Used in \\
\midrule
LOSO & one fit per source on all other sources; the held-out source is scored (\cref{alg:buildX}) & every development analysis (v1--v12) \\
Nested LOSO & as LOSO, plus inner LOSO selection within each training set (\cref{sec:nested}) & v4a, v7c2, v8d \\
Frozen external & one fit on S1; all external sources scored by the same model, with S1 imputation medians & v6a (set E), v8b (set X) \\
Frozen out-of-source & models fitted on the previous NICER data before the new data are read; each source's new observations are scored by the fit without that source & v11 (reference v10), v13 (reference v12) \\
All-data model & one fit on all reference data, for sources never used before & v11 (never-used sources) \\
\bottomrule
\end{tabularx}
\end{table}
